\documentclass[11pt]{letter}

\usepackage[utf8]{inputenc}
\usepackage[T1]{fontenc}
\usepackage[margin=1in]{geometry}
\usepackage{hyperref}
\usepackage{siunitx}
\urlstyle{same}

\sisetup{per-mode=symbol, group-separator={,}}

\signature{Jonathan Mace\\Microsoft Research}
\address{Jonathan Mace\\Microsoft Research\\Redmond, WA\\jonathanmace@microsoft.com}
\date{\today}

\begin{document}

\begin{letter}{Editor\\Journal of Sound and Vibration\\Elsevier}

\opening{Dear Editor,}

\small

Please consider our manuscript, ``Geometry-dependent identifiability in spectral
inverse problems: oblate lifting without prolate improvement,'' for publication
in the \emph{Journal of Sound and Vibration} as a Short Communication.

The paper addresses a precise question in spectral inverse problems for
fluid-loaded elastic shells: does breaking spherical symmetry always improve
parameter recoverability?  Using a five-mode flexural frequency map and the
scaled Jacobian formalism, we show that the answer depends on the
\emph{direction} of the symmetry breaking.  The equivalent-sphere reduction is
rank deficient ($\sigma_3=0$), and this singularity persists at the exact
sphere.  Along the oblate branch, the weakest singular value grows
quadratically, $\sigma_3(\varepsilon)=\lambda_1\varepsilon^2+O(\varepsilon^4)$,
reducing the condition number from $\kappa=244$ (Ritz sphere) to
$\kappa_\mathrm{oblate}=\num{69.4}$ at the canonical geometry.  The prolate
continuation, by contrast, yields $\kappa_\mathrm{prolate}=156$ at $c/a=1.5$
with no systematic improvement across the tested range.

This result is concise but non-trivial: it shows that ``symmetry breaking
helps'' is false as a general rule for spectral inverse problems, and that the
identifiability gain is tied to how curvature redistribution separates the
geometric sensitivity directions.  The finding has direct implications for
inverse problems in biomedical acoustics (where abdominal organs are
characteristically oblate) and non-destructive evaluation of thin-walled
vessels.

The manuscript is part of a broader research programme on the vibration
mechanics of fluid-filled shells.  The companion Paper~1 (submitted to JSV)
establishes the canonical oblate shell model and its flexural spectrum, while
Paper~8 (submitted to Inverse Problems) provides the full identifiability
analysis including Newton inversion and Cram\'er--Rao bounds.  The present Short
Communication distils the oblate-versus-prolate asymmetry result into a
self-contained form suitable for the JSV readership.

We believe the manuscript fits JSV because the journal regularly publishes work
on structural dynamics, inverse problems, and fluid--structure interaction, and
the Short Communication format is appropriate for the focused nature of this
contribution.

The manuscript is original, has not been published previously, and is not under
consideration elsewhere.  Both authors have approved the submission.  Brian R.\
Mace's contact email is \texttt{b.mace@auckland.ac.nz}.

\closing{Yours sincerely,}

\end{letter}

\end{document}
